\chapter[Completed: Super Odometry2.0: Hierarchical Adaptation Enables Robust Odometry Towards
All-degraded Environments]{Super Odometry2.0: Hierarchical Adaptation Enables Robust Odometry Towards
All-degraded Environments}

\label{chapter:superodom2.0}
\fancyhead[LO]{\makebox[\textwidth][r]{Tartan IMU: A Light Foundation Model for Inertial Positioning in Robotics}}

\begin{framed} The material from this chapter is taken from the following paper. \\
Zhao, Shibo, et al. Hierarchical Adaptation Enables Robust Odometry Towards
All-degraded Environments  Under review 2025.
\end{framed}




\section{Introduction}
Odometry is an important technique to estimate the position and orientation of robots over time, while also allowing for the 3D geometry reconstruction of surrounding environments. It plays a crucial role in robotics, enabling spatial awareness and serving as a foundation for both high-level tasks such as navigation and exploration, and low-level functions like path planning and control~\cite{thrun2002probabilistic}. As a result, odometry systems are widely used in robotic applications including off-road driving~\cite{selfdriving}, search-and-rescue~\cite{Scherer:2022}, and planetary exploration~\cite{mahlknecht2022exploring}. 

However, maintaining robust and resilient odometry performance in real-world applications remains challenging due to a range of environmental degradations, such as poor lighting, geometrical degradation, motion blur from aggressive maneuvers, and extreme weather. These degradation cases make odometry system hard to extract and track reliable features from sensors over time, often leading to incorrect data association\cite{cadena2016past}.

% complicate the extraction of reliable features from sensors and their consistent tracking over time, often leading to unstable pose estimation. % 
% incorrect data association.


To overcome these environmental degradations, researchers have explored multi-modal sensor fusion methods, implemented as loosely-coupled \cite{shen,wan2001unscented,zhang2014loam,zhang2015visual,complementaryodometry,camurri2020pronto} or tightly-coupled \cite{graeter2018limo,zuo2019lic,wisth2021unified,r3live} frameworks. Although these approaches have substantially advanced the field,
none of them can fully handle all forms of environmental degradation, including visual, geometrical, and their combined degradation ~\cite{zhao2024subt}. 
This shortcoming continues to hinder the development of robust and resilient autonomous navigation systems capable of long-term operation in any location and at any time.
The key challenges can be summarized as follows: 

% \textit{Generalizing across Diverse Scenes:}
% Most methods need a significant amount of tuning for different types of environments. For example, indoor and outdoor scenes may need a very different set of parameters. In practice, it is very important to be able to adjust to various types of scenes and adapt to environmental changes without manual tuning.

\textit{Adaptation Across Varied Degradation Scenarios:}
Most methods need a significant amount of tuning for different types of environments since they are built on a "rigid" framework with a "fixed" parameter setting~\cite{cadena2016past}. For example, a method optimized for outdoor conditions may perform poorly in indoor environments such as caves.  These frameworks are difficult to adapt to varying types of degradation.

% are built on “rigid” frameworks that passively integrate sensor measurements or add additional sensors to mitigate data degradation.


\textit{Being Robust to Severe and Persistent Sensor Degradation:}
Existing SLAM solutions encounter significant challenges when multiple sensors experience degradation or remain off-nominal over extended periods. For example, persistent smoke in wildfire environments can severely impact both visual and LiDAR sensors, leading to substantial navigation difficulties~\cite{ebadi2023present}. 

\textit{Balancing Robustness and Efficiency:}
Most approaches enhance robustness by integrating data from all available sensors to ensure redundancy and safety. However, this approach often leads to computationally expensive optimization problems~\cite{cadena2016past}. Achieving a balance between robustness and efficiency remains a significant challenge.

% \textit{Balancing Redundancy and Computational Efficiency:} most approaches fuse all available sensor data to enhance robustness, which can lead to excessive data redundancy and high computational demands. A balance between robustness and lightweight computation is difficult.

% \textit{Being Fail-Safe and Failure-Aware:} An ideal SLAM solution should be \textit{fail-safe} and \textit{failure-aware} (i.e., the system needs to be aware of imminent failure and provide recovery mechanisms)\cite{cadena2016past}. However, this capacity is lacking in most existing odometry systems.

% \textit{Robust Performance in Persistent Sensor Degradation:}
% Current SLAM solutions struggle to maintain robust performance when certain sensors are compromised over extended periods, such as in environments with persistent smoke or sensor-drop scenarios.

% These limitations significantly restrict the range of environments where robots would have outsize value like search-and-rescue missions in wildfires, earthquakes, sandstorms, and snow.




% To address these challenges, we believe a robust odometry system should adapt or adjust its operation based on the severity of degradation. For example, mild degradation(e.g., low-lighting conditions), might only require minor adjustments to maintain performance. In contrast, severe degradation( e.g., obstructed vision from dense smoke or heavy dust), requires more sophisticated strategies~\cite{zuo2019lic, wisth2021unified, r3live}. However, many current algorithms treat all types of degradation uniformly. As a result, they struggle to deliver
% optimal performance in challenging environments where an adaptive response to increasing degradation is essential for maintaining reliable operation~\cite{cadena2016past}\cite{ebadi2023present}. 

% For example, these methods might be difficult to work well in supporting various environments and different robot platforms to achieve year-long persistence~\cite{cadena2016past}\cite{ebadi2023present}.





\begin{figure*}[h]
    \centering
    \includegraphics[width=0.98\textwidth, keepaspectratio]{figs/ch7/science_robotics/concept.pdf}
    \caption{\textbf{Mechanism of sensory adaptation in humans and robots. } \textbf{(A)} humans evolved to use a hierarchical sensory adaptation mechanism to estimate ego-motion effectively even in challenging environments. Strategies range from basic adjustments (e.g. visual self-attention) to complex multi-sensory integration and cognitive processes (e.g., path integration). \textbf{(B)} Inspired by the above, we propose a hierarchical adaptation strategy that initiates with adaptive feature selection for visual degradation, moves to an adaptive state direction for LiDAR degradation, applies adaptive factor (engine) selection for mixed degradation, and employs learning-based  inertial odometry during complete degradation. Image credits for brain, eye, vestibular, and depth perception anatomy: \cite{duckgogit, AnhatomyStuff,oktar2019convolutional,NEUROSCIENTIFICALLY}.}
    \label{fig:concept}
\end{figure*}

In contrast, humans and animals exhibit exceptional adaptability when navigating complex and ever-changing environments, relying on multiple sensory inputs~\cite{ptito2006sensory}.
% when it comes to multi-modal sensing, allowing them to estimate ego-motion effectively, even in complex and ever-changing environments~\cite{ptito2006sensory}. 
% They achieve this by utilizing a combination of external cues, such as visual input from the eyes and tactile sense from the hands, along with internal cues from the vestibular system. 
% For instance, as people navigate a poorly lit corridor and eventually enter a smoke-filled room with conditions deteriorating from mild to extreme (Fig.~\ref{fig:concept} A). 
For instance, as individuals move through a dimly lit corridor and enter a smoke-filled room with conditions progressively worsening from mild to extreme, they employ various strategies to maintain spatial awareness, as shown in Fig.~\ref{fig:concept} A.
These approaches range from basic adjustments, such as selective attention to key visual cues \cite{nuthmann2020salience}, to more sophisticated multi-sensory integration involving vestibular and visual inputs~\cite{ptito2006sensory, braunstein2014depth}. In dense smoke, where vision is fully obscured, the human brain shifts reliance from visual input to the proprioceptive sensory and vestibular system to support path integration~\cite{bostelmann2020path}.
Together, these mechanisms create a multi-layered, adaptive response system that allows humans to navigate effectively even in sensory-degraded environments.
This hierarchical framework, from basic sensory tuning, multi-sensory integration, and reliance on alternative sensory cues from brains, provides a compelling model for the design of next-generation robotic odometry systems.
% We observe people adapt in multiple ways to localize themself in such challenging environments. 
% Their adaptation strategies range from basic adjustments (e.g., selective attention on key visual cues~\cite{nuthmann2020salience}) to complex multi-sensory integration (e.g., utilizing both vestibular and visual cues~\cite{ptito2006sensory, braunstein2014depth}). 
% In dense smoke, where vision is fully obscured, the human brain shifts reliance from visual input to the proprioceptive sensory and vestibular system to support path integration~\cite{bostelmann2020path}.
% Together, these mechanisms create a multi-layered, adaptive response system that allows humans to navigate effectively even in sensory-degraded environments. This hierarchy of basic sensory adjustment, multi-sensory integration, and reliance on alternative cues from brains offers an insightful model for designing robotic odometry systems. 

% We believe this adaptive hierarchy could enable more effective spatial awareness across a range of difficult environments.

% and cognitive processes (e.g., path integration only from vestibular cues~\cite{bostelmann2020path}), shown in Fig.~\ref{fig:concept} A (a-d).
% In low-light conditions, humans adjust pupil size and focus on key visual cues like illuminated signs~\cite{nuthmann2020salience}. In featureless corridors, where depth perception is unreliable, they rely on nearby landmarks, such as doors, to estimate their position~\cite{ptito2006sensory, braunstein2014depth}. In dense smoke, 
% They leverage their path integration,  proprioception, vestibular cues, and spatial memory to approximate location without external visual feedback.
% \textit{This adaptive hierarchy enables effective navigation and spatial awareness across a spectrum of challenging environments.} 

Inspired by these biological strategies, we suggest that a robotics odometry system should adopt a hierarchical adaptive strategy ranging from light adjustments for mild degradation to more computationally intensive adjustments for severe degradation. This multi-level adaptation scheme can achieve a balance between data redundancy for robustness and computational efficiency.
Moreover, the odometry system should have robust path integration capabilities like those of humans, enabling it to function effectively in perceptually degraded environments. Beyond human-level performance, the system should also be generalizable to various robotic platforms and diverse environments. By processing more and more data, the odometry system can progressively enhance its performance.


To achieve the above, we present a novel sensor fusion framework called Super Odometry2.0, which achieves robust performance in all-degraded environments. It is empowered by two main techniques: 1) a hierarchical adaptive framework that seamlessly integrates various adaptation strategies, enabling robust performance across varying levels of environmental degradation. 2) a learning-based inertial odometry with exceptionally low drift, capable of online adaptation to new environments and platforms.  
% emulates the adaptation mechanisms of human multimodal sensing. 
Our major contributions are as follows:




% We classified the sensor degradation with four distinct types of sensory degradation: visual degradation, depth (geometric) degradation, mixed degradation, and complete degradation. 

% \paragraph{Visual Perception Degradation}
% In low-light conditions (Fig. \ref{fig:concept} B-a), we proposed adaptive feature selection. Our key observation is that not all visual or geometric features contribute equally to state estimation. Therefore, we prioritize the most valuable features for motion estimation. 

% \paragraph{Geometric Perception Degradation}
% When robots navigate through featureless long corridors (Fig. \ref{fig:concept} B-b), we propose adaptive state direction, which  can not only predict alignment risks, and
% estimate the observability of state direction, but also can actively incorporate
% pose priors from other odometry sources before failure occurs. 

% \paragraph{Mixed Perception Degradation}
% In more challenging environments (Fig. \ref{fig:concept} B-c) where both vision and depth perception are compromised, we propose an adaptive engine (factor) selection approach to address long-term sensor-level degradation.


% \paragraph{Complete Perception Degradation}
% In extreme scenarios, such as dense smoke (Fig. \ref{fig:concept} B-d), visual cues become entirely unreliable. In these "blind" situations, we introduce a novel deep learning-based inertial odometry designed to
% capture and learn the motion patterns of robotics since it is not affected by the environment. 
% This hierarchical adaptation process demonstrates how humans and animals have evolved to respond to varying degrees of environmental degradation. 
% Strategies range from basic adjustments (e.g., pupil dilation) to complex multi-sensory integration and cognitive processes (e.g., path integration). \textit{This adaptive hierarchy enables effective navigation and spatial awareness across a spectrum of challenging environments.} Inspired by these biological strategies, we propose a novel hierarchical adaptation framework for odometry systems, called Super Odometry, which emulates the adaptation mechanisms of human multimodal sensing. 
% This adaptation framework enables our system to:
% a) Scale its response based on the severity of environmental degradation b) Efficiently allocate computational resources c) Integrate multiple sensory inputs adaptively d) Maintain robustness through redundancy and fault tolerance.

% We group adaptive feature
% selection, state direction, and engine selection under the term adaptive
% sensor fusion, as they are based on classic factor graph optimization.
% We treat learning-based inertial odometry separately, as it relies on a
% deep learning network. Together, these components form what we refer
% to as a hierarchical adaptation framework.
% In summary, our major contribution is as follows:

\begin{figure*}[p]
   \centering
    \includegraphics[width=1.0\textwidth,height=0.88\textheight,keepaspectratio]{figs/ch7/science_robotics/various_degraded.pdf}

    \caption{\textbf{Robust odometry across diverse conditions with various sensors and robots.}  We tested our method under a wide range of challenges, which covers various scenarios with visual degradation (A-D), geometric degradation (E-H), mixed degradation(I-L), complete degradation (M-P) and other scenarios such as caves, off-road, steep slopes, and featureless deserts (Q-T).} 
    \label{fig:various_scenes}
\end{figure*}


\textbf{Hierarchical Adaptation:} 
We propose the first hierarchical adaptive framework to address degradation along a spectrum, from mild to severe, and scale its responses accordingly to maintain efficiency and robustness. 
The system is dynamically reconfigurable, functioning as a multi-level scheme and each level focuses on different types of degradation. 
Specifically, the system initiates with adaptive feature selection to mitigate mild visual degradation, then transitions to adaptive state direction for moderate geometric degradation, proceeds to adaptive factor (engine) selection for mixed degradation,  and ultimately relies on a learning-based inertial odometry module for complete degradation, as illustrated in \fref{fig:concept} B.

Lower levels provides rapid and resource-efficient adaptations to address less severe issues. If these measures are insufficient, higher levels provide more complex and computationally intensive interventions to support state estimation recovery. By structuring adaptation hierarchically, our approach significantly enhances the resilience of robotic navigation systems, ensuring reliable operation under nearly all challenging conditions.

% We propose a hierarchical adaptation framework to address degradation along a spectrum, from mild to severe, and scale its responses accordingly to maintain efficiency and robustness.  The system is dynamically reconfigurable, functioning as a multi-level scheme and each level focuses on different types of degradation. Specifically, the system initiates with adaptive feature selection for visual degradation, moves to an adaptive state direction for geometric degradation, applies adaptive factor selection for mixed degradation, and employs learning-based inertial odometry for complete degradation (\fref{fig:concept} B). Lower levels handle basic, fast, and resource-efficient adaptation and higher levels handle complex and computationally demanding processes.
% If low-level adaptation mechanisms fail, high-level mechanisms can intervene to support recovery.  This hierarchical method enhances the resilience of state estimation, allowing the robot to effectively navigate through all forms of challenges.


\textbf{Heterogeneous Learning-based Inertial Odometry:} 
Humans do path integration using proprioceptive sensors and vestibular systems in severely degraded environments. Inspired by this, we introduce a heterogeneous learning-based inertial odometry that captures different robot motion dynamics and predicts 3D poses in real-time. Our model is trained on \textit{hundreds of hours} of IMU data from diverse robotic platforms, enabling strong generalization across different robot platforms. This comprehensive training approach not only outperforms existing specialized models tailored to specific robots but also mitigates the risks of overfitting. 
Experimental results indicate that the proposed method can maintain accurate pose estimates with minimal drift over several minutes, even in previously unseen environments.
% Results show that the system can maintain accurate pose integration for several minutes with minimal drift, even in unseen environments.


\textbf{Learning and Adaptation Over Time:}
We developed an online adaptation algorithm for the data-driven inertial odometry using self-supervised learning\cite{wang2024imperative}.
This adaptation system is formulated as a bilevel optimization: the upper-level is an inertial odometry network optimization, which is subjected to a lower-level factor graph optimization~\cite{dellaert2017factor}.
% By integrating inertial odometry into the system we leverage the outputs from factor graph sensor fusion~\cite{dellaert2017factor} as supervision signals. 
% This allows the inertial odometry to learn underlying motion patterns in real-time without requiring any ground-truth.
The factor graph provides free pose labels, enabling the inertial odometry to learn motion patterns and predict poses with minimal drift.
In return, the inertial odometry offers a motion prior for the factor graph, serving as a robust fallback for state estimation when other sensors fail.
% In return, the inertial odometry offers a motion prior to the factor graph, serving as a fallback for state estimation when other sensors fail. As data-driven inertial odometry processes more motion patterns, its performance continually improves. 
For the first time, we elevate IMUs to an independent role in sensor fusion, positioning them as equally important as camera and LiDAR.

% As a result, our learning inertial odometry can adapt to new environments and new robots autonomously without manual tuning. When other sensors, such as visual or LiDAR, are functioning well, the inertial odometry receives feedback and fine-tunes itself. When other sensors are compromised, the inertial odometry takes over and estimates the state accurately for as long as a few minutes. This approach mirrors how humans adapt by using both external visual cues and internal signals from the vestibular system. 
% We demonstrate that odometry systems could benefit from the inertial-sensor-only method, marking the first instance where IMU plays a central role as important as visual and LiDAR sensors.

\textbf{Extensive Evaluation:} We extensively test the Super Odometry2.0 across a wide range of robotic platforms, including aerial, wheeled, legged, and handheld robots, to accommodate various locomotion modes as shown in \fref{fig:various_scenes}. In environments with different levels of degradation, it consistently outperformed state-of-the-art techniques in terms of robustness, accuracy, and efficiency. We also show that Super Odometry2.0 has exceptional resilience to environmental degradation, achieving \textit{near-zero} failure over a wide variety of trajectories, accumulating more than \textbf{200} kilometers and \textbf{800} operational hours in six years.

We believe Super Odometry2.0 represents a significant step towards a resilient and robust sensor fusion solution capable of operating reliably anytime, anywhere, and in any degraded environment. We expect this advancement will enhance long-term robotic autonomy and ensure the safety of robots.  For more details on motivation and methods, please refer to Supplementary \textbf{Movie S1}.


\section{Related Works}
\label{chapter7:related_works}

Super Odometry integrates an \textit{adaptive sensor fusion strategy with a learning-based inertial odometry} in a self-supervised manner. The adaptive sensor fusion framework assesses degeneracy from low to high levels and provides supervision signals to the learning-based inertial state estimator. In return, the learning-based inertial state estimator learns the motion dynamics of robots and offers motion priors to the sensor fusion framework. In addition, it allows the system to run without failure for up to a few minutes in extreme degradation where all exteroceptive sensors fail.   
In the following paragraphs, we will give an overview of the \textbf{adaptive sensor fusion} - which consists of adaptive feature selection, adaptive state direction, adaptive engine selection, and the learning-based inertial odometry.  
% We group adaptive feature selection, state direction, and engine selection under \textbf{adaptive sensor fusion}, as they are based on classic factor graph optimization. We treat \textbf{learning-based inertial odometry} separately, as it relies on a deep learning network. Together, these components form a \textbf{hierarchical adaptation framework}. 


% The system design follows a key observation: an inertial measurement unit (IMU) is lightweight, computationally efficient, and unaffected by perceptual degradation. Therefore, our approach elevated IMU to the same level of importance as visual and LiDAR sensors.

\subsection{Adaptive Sensor Fusion}  

Handling multiple degradations is a major challenge for odometry systems. 
Although several strategies have been proposed to alleviate this issue, e.g., degeneracy mitigation \cite{tagliabue2021lion,tuna2023x}, uncertainty estimation\cite{talbot2023principled,brossard2020new} and observability analysis \cite{degeneracy,zhen2019estimating,ding2021degeneration,hinduja2019degeneracy,tagliabue2021lion}, none of these can overcome all types of degraded environments.
% This is because existing methods require prior knowledge about the environment or heavily rely on manual parametric tuning.
% Most methods perform observability analysis in the back-end, which can be too late to prevent failures due to insufficient data association \cite{degeneracy}.
To solve this, we introduce adaptive sensor fusion, including adaptive feature selection, adaptive state direction, and adaptive engine selection.

\textit{\textbf{Adaptive Feature Selection:}~}
In visually degraded environments, we observed that not all visual features contribute equally to state estimation (Fig.~\ref{fig:concept} B). Therefore, we prioritize the most valuable features by analyzing the uncertainty associated with each feature~\cite{muhle2022probabilistic}.
% For cameras, we identify valuable feature points based on appearance and analyze their contribution to motion estimation. For LiDAR, we assess the observability of geometric features and estimate the covariance matrix of the scan point cloud, allowing us to predict the uncertainty and quality of ICP-based scan matching without heavy optimization.


\textit{\textbf{Adaptive State Direction:} ~}
In geometrical degradation, some directions, such as along a long and windowless corridor, cannot be well constrained by the LiDAR odometry shown in Fig.~\ref{fig:concept} B.
Our method can identify these poorly constrained directions and actively incorporate pose priors from other odometry sources to enhance robustness.

\textit{\textbf{Adaptive Engine Selection:} ~}
To address mixed degradation from visual and geometric sources, we further employ factor or engine level uncertainty analysis to actively switch auxiliary sensors based on degeneracy type (Fig.~\ref{fig:concept} B).



\subsection{Heterogeneous Learning-based Inertial Odometry} 
% Describe Learning IMU odometry new direction 
% Describe how successful of learning IMU odometry 
% Describe their limitation of learnging IMU odometry 
% Describe what is our strenth for these limitation 
% (Adaptation Capability for quick adapt, Genralize Capability for any robot, )
The emergence of deep learning offers new possibilities for extracting information from IMU data and developing motion models for robots. 
Several studies such as IONet \cite{chen2018ionet}, TLIO \cite{liu2020tlio}, RIDI \cite{yan2018ridi}, and RoNIN \cite{chen2021rnin} have demonstrated that learning-based IMU estimators are promising in position estimation and can be used in complete degraded environments such as smoke as shown in Fig.~\ref{fig:concept} B. However, most of these methods are trained on one dataset for one specific platform. They do not generalize across datasets or platforms~\cite{buchanan2022deep}. To solve this, we made two efforts. 

\textit{\textbf{Pre-trained IMU Model on Large Datasets:}}
we propose a \textit{general} IMU model, featuring a shared backbone that enables cross-platform motion estimation. It
is trained on large-scale, high-quality, and diverse IMU data \cite{zhao2024subt} \cite{sivaprakasam2024tartandrive} comprising hundreds of hours of data from various platforms, including drones, quadrupeds, cars, and humans, in a variety of scenarios, including cave exploration, drone racing, and indoor navigation. This broad training enables our \textit{general} model to outperform any specific models~\cite{chen2018ionet,liu2020learning,chen2021rnin,yan2018ridi} and to be applied to various robot platforms with different motion patterns.

\textit{\textbf{Online Adaptation to Unseen Datasets:}} Although our model is trained on a large dataset, there are always out-of-distribution motion patterns that are not covered. 
% no guarantee that it can reliably predict poses for every unseen robot platform or motion pattern. 
Furthermore, collecting sufficient data to fine-tune the pre-trained IMU network during deployment is not always feasible. This requires the network to adapt itself \textit{seamlessly} to changing environments and platforms. To address this challenge, we propose an online adaptation scheme for our pre-trained IMU model, enabling it to adjust to novel environments continuously in a self-supervised manner.
Unlike conventional learning approaches where the pre-trained model remains static during inference, our method eliminates the boundary between the training and testing phases — \textbf{we learn as we operate}. To enable fast adaptation, we integrate Low-Rank Adaptation (LoRA) \cite{hu2021lora} into our online adaptation process, significantly boosting real-time performance. Our inertial odometry model achieves 200 FPS with real-time refinement.

To the best of our knowledge, this is the \textbf{first} IMU-based estimator that demonstrates such a high level of generalization across diverse robot platforms, fast adaptation to unseen environments, and real-time predictions with minimal drift.

\begin{figure*}[p]
   \centering
\includegraphics[width=1.0\textwidth,height=0.85\textheight,keepaspectratio]{figs/ch7/science_robotics/Method2.pdf}
     \caption{\textbf{System Overview.} \textbf{(A)} A hierarchical adaptation framework that integrates adaptive sensor fusion with learning-based IMU odometry. The sensor fusion provides free pose labels to guide learning-based IMU odometry reduing drift, while the IMU odometry supplies motion priors and serves as a fail-safe. \textbf{(B)} Details for the learning-based IMU odometry module. \textbf{(C)} Process of adaptive feature selection, adaptive state direction estimation, and adaptive engine selection.}
    \label{fig:method}
\end{figure*}




\section{Problem Statement}
\label{chapter7:problem_statement}

\paragraph{Term Clarification} Before outlining our methodology, we first clarify the key technical terms used in this paper. We propose hierarchical multi-level mechanisms to address environmental degradation, ranging from mild to extreme. These mechanisms include adaptive feature selection, adaptive state direction, adaptive engine selection, and learning-based inertial odometry. For clarity, we group adaptive feature selection, state direction, and engine selection under the term \textbf{adaptive sensor fusion}, as they are based on classic factor graph optimization. We treat \textbf{learning-based inertial odometry} separately, as it relies on a deep learning network. Together, these components form what we define as a \textbf{hierarchical adaptation framework}. 

\section{Approach}
The key aspect of our method is a hierarchical adaptation framework that integrates adaptive sensor fusion with learning-based IMU odometry in a \textbf{self-supervised manner}, as shown in Fig.~\ref{fig:method} A. This framework operates as a bilevel optimization\cite{wang2024imperative}, where the adaptive sensor fusion acts as the "teacher" model and the learning-based IMU odometry acts as the "student" model. 
When visual or LiDAR sensors are functioning well, the learning-based IMU odometry receives feedback from adaptive sensor fusion and learn the the motion pattern of robots. When camera or LiDAR sensors are compromised, the learning IMU odometry takes over in the sensor fusion pipeline, reliably estimating the robot state.
Together, this cooperative symbiosis significantly enhance the robustness, resilience, and adaptability of odometry systems in challenging environments.
% Because of space constraints, detailed algorithms, and implementation specifics for handling various types of degeneracy, we refer readers to the Supplementary Materials. 

\subsection{Adaptive Sensor Fusion} 
% We believe that robots should possess adaptive capabilities similar to those of animals, enabling them to perceive their surroundings and localize themselves robustly in various degraded environments. To achieve this,
We propose adaptive sensor fusion, which hierarchically selects the most meaningful information from  parameters, features, states, and engines, ranging from low-level to high-level. It consists of adaptive feature selection, adaptive state selection, adaptive engine selection, and adaptive parameter tuning.


\textbf{Adaptive Feature Selection~}
 As shown in Fig.~\ref{fig:method} C-Step1, our key observation is that not all the features, whether visual or geometric features in nature, contribute equally to the accuracy and robustness of state estimation. Therefore, it is important to select the most valuable features for motion estimation. For the camera, we analyze the Hessian matrix of KLT tracking \cite{hwangbo2009inertial} and determine the distribution of intensity residuals by calculating their covariance matrix $\boldsymbol{\Sigma}_{\mathrm{SE(2)}}$ using small patch sizes \cite{muhle2022probabilistic}.
Taking into account the estimated rotation $\boldsymbol{R}_\theta \in \mathrm{SO}(2)$ from the odomery system, we can then compute the rotated covariance for each individual 2D feature patch.

\begin{equation}
\boldsymbol{\Sigma}_{2 \mathrm{D}, \theta}=\boldsymbol{R}_\theta \boldsymbol{\Sigma}_{\mathrm{SE(2)}} \boldsymbol{R}_\theta^{\top}    
\end{equation}


% \begin{equation}
% \begin{cases}
% 	\mathbf{e}_{i}^{\mathrm{po} \rightarrow \mathrm{po}} = \mathbf{q}_{i} - (\mathbf{R}\mathbf{p}_{i} + \mathbf{t}) & \text{for } q_{i}, p_{i} \in F^{po}_{i+1} \\
% 	\mathbf{e}_{j}^{\mathrm{po} \rightarrow \mathrm{li}} = \mathbf{v}_{j} \times \left(\mathbf{q}_{j} - (\mathbf{Rp}_{j} + \mathbf{t})\right) & \text{for } q_{j}, p_{j} \in F^{li}_{i+1} \\
% 	\mathbf{e}_{k}^{\mathrm{po} \rightarrow \mathrm{pl}} = \mathbf{n}_{k}^{\top} \left(\mathbf{q}_{k} - (\mathbf{Rp}_{k} + \mathbf{t})\right) & \text{for } q_{k}, p_{k} \in F^{pl}_{i+1}
% \end{cases}
% \end{equation}

% where \( \mathbf{e}_{i}^{\mathrm{po} \rightarrow \mathrm{po}} \), \( \mathbf{e}_{j}^{\mathrm{po} \rightarrow \mathrm{li}} \), and \( \mathbf{e}_{k}^{\mathrm{po} \rightarrow \mathrm{pl}} \) represent the point-to-point (line, plane) distances. Here, \( \mathbf{v}_{j} \) and \( \mathbf{n}_{k} \) are the principal and normal directions of the respective features, and \( \mathbf{T}_{i+1} = \{\mathbf{R}, \mathbf{t}\} \) is the optimal transformation we aim to estimate.


For LiDAR, we evaluate the quality of correspondences based on the fit of neighboring points to point, line, or plane distributions \cite{demantke2011dimensionality}:
\begin{equation}
 w_i^{\mathrm{po} \rightarrow \mathrm{li}}= \frac{\sigma_1 - \sigma_2}{\mu}, \quad  w_i^{\mathrm{po} \rightarrow \mathrm{pl}} = \frac{\sigma_2 - \sigma_3}{\mu}, \quad  w_i^{\mathrm{po} \rightarrow \mathrm{po}} = \frac{\sigma_3}{\mu},    
\end{equation}

Where \( \sigma_i \) are the eigenvalues of the Principal Component Analysis (PCA) used for point-to-point, point-to-line, point-to-plane distance error metric\cite{zhao2021super}. The normalization coefficient \( \mu = \sum_{i=1}^3 \sigma_i \) is used and
\( W_l = \{ w_i^{\mathrm{po} \rightarrow \mathrm{po}}, w_i^{\mathrm{po} \rightarrow \mathrm{li}}, w_i^{\mathrm{po} \rightarrow \mathrm{pl}} \} \in [0, 1] \) describe the linearity (\( w_i^{\mathrm{po} \rightarrow \mathrm{li}} \)), planarity (\( w_i^{\mathrm{po} \rightarrow \mathrm{pl}} \)), and scatter (\( w_i^{\mathrm{po} \rightarrow \mathrm{po}} \)) of the features.



% If \( \sigma_1 \gg \sigma_2, \sigma_3 \simeq 0 \), \( a_{1D} \) will be greater than the two others so that the dimensionality labelling \( d^*(\mathcal{V}^r_P) \) results in 1. Contrariwise, if \( \sigma_1, \sigma_2 \gg \sigma_3 \simeq 0 \), \( a_{2D} \) i.e., the planar behavior will prevail. At last, \( \sigma_1 \simeq \sigma_2 \simeq \sigma_3 \) implies \( d^*(\mathcal{V}^r_P) = 3 \).

% We chose \( \mu = \sigma_1 \) because then: \( a_{1D} + a_{2D} + a_{3D} = 1 \), so that the three features can be considered as the probabilities of each point to be labelled as 1D, 2D, or 3D. It will help us to select the most appropriate neighborhood size by finding which radius favors the most one dimensionality (see Equation 3).

% In this equation, \( p_i \) are the neighboring points of the extracted features, and \( \bar{p} \) is their mean. For point-to-line (or point-to-plane) distance errors, we use the normal direction \( \mathbf{n}_k \) of the features and define \( A = I - \mathbf{n}_k^{\top} \mathbf{n}_k \) for point-to-line and \( A = \mathbf{n}_k^{\top} \mathbf{n}_k \) for point-to-plane. \( d_{max} \) represents the maximum distance error for point-to-line (or point-to-plane) correspondences, and \( k \) is the number of neighboring points.

This approach enables the selection of the most meaningful features while reducing the computational load.

%For more details, we refer readers to our supplementary material.



% various and complex environments, Super Odometry also achieve self-tuning capability to adapt all-terrain environments. To improve efficiency, Super Odometry adopts adaptive voxel size for different indoor and outdoor environments. To improve robustness in subterranean environments, Super Odometry enable or disable dust filter based on the quality of geometric features. To overcome different terrain environments, Super Odometry adopts multi-hypothesis for orientation estimation.

% Deploying efficient algorithms on heterogeneous platforms in subterranean environments remains an open challenge due to the limitations of onboard computation needed for autonomous operation. To solve this problem, self-tuning capabilities is very important for the SLAM system to adapt the selection of the system parameters to the different scenario. Matteo Palieri et al \cite{palieri2020locus} present LOCUS architecture which can be adapted to heterogeneous robotic platforms with diverse sensor inputs and computational capabilities. Then, Andrzej et al \cite{reinke2022locus} extend LOCUS and include an adaptive voxel grid filter that maintains the desired computation load regardless of the environment’s geometry. 
% Milad Ramezani et al proposed Wildcat \cite{ramezani2022wildcat} which performed continuous-time trajectory representation with an efficient pose-graph optimization module that seamlessly supports heterogeneous platforms.

\textbf{Adaptive States Selection~}
% 1. multi sensor fusion has different strength on state direction esttimation base on the change of environemtns. 
% 2. One great example, in long orridor environment, lidar odometry fail in forward but good extimation in side direction.
% However, most of methods didn't consider the state direction and direclty fuse all the direction from multiple sensor and regard the fusion equally. 
% 3 In constrast, our method analyze the observability in advance and 
Various forms of degradation hinder the robustness of existing odometry systems \cite{xu2021fast}\cite{LVI-SAM}. 
% However, most current approaches only passively incorporate sensor information and are unable to identify degenerate directions in the estimated
% state\cite{degeneracy,zhen2019estimating,ding2021degeneration,hinduja2019degeneracy}.
However, most approaches require prior knowledge of the environment to determine well-conditioned directions of state estimation and cannot adapt to changing environments\cite{degeneracy,zhen2019estimating,ding2021degeneration,hinduja2019degeneracy}.
% Zhang et al. \cite{degeneracy} analyzed the information matrix of visual/LiDAR odometry and selectively fused data only in well-conditioned directions. This technique has been further applied to tasks such as localization \cite{zhen2019estimating, ding2021degeneration} and scan registration \cite{hinduja2019degeneracy}. However, these methods require prior knowledge of the environment to determine well-conditioned directions and cannot adapt to changing environments. 
Moreover, some methods perform degeneracy analysis \cite{ebadi2021dare} and robust kernel function \cite{chebrolu2021adaptive} during or after the optimization, which can be too late to predict potential failures \cite{degeneracy} and insufficient data association already causes unreliable optimization. 

To address these issues, we perform a degeneracy analysis by calculating the observability at the front end of LiDAR odometry to identify ill-conditioned directions and mitigate the potential failure before optimization shown in Fig.~\ref{fig:method} C-Step2.  We focus on LiDAR odometry for degeneracy analysis rather than visual odometry because LiDAR provides more reliable observability and uncertainty measurements for each state direction. Once LiDAR determines the uncertainty in each state direction, the visual sensor can complement it by adjusting the corresponding relative weights. To achieve this, we proposed to predict the alignment risk of ICP scan registration and then calculate the observability $\mathbf{Obs}_i$ for each point in every state direction. \textit{Then, we can assign labels to each correspondence based on the state direction it primarily influences. } Because of space constraints, please see Supplementary Materials on \textbf{alignment risk prediction and observability} for each point-to-plane constraint. 

To comprehensively evaluate the observability of an entire scan, we develop a statistical approach that analyzes the distribution of observability labels across different state directions. Our method normalizes these labels to a [0-1] range, creating a confidence metric that quantifies the reliability of state estimation. We compute normalized observability metrics for both translational and rotational components using the following formulations:

 
\begin{align}
\mathbf{Cov}_{trans} &= 3\left[\frac{N_{\mathbf{Obs}_x} }{N_{\mathbf{Obs}_{sum}} },\frac{N_{\mathbf{Obs}_y} }{N_{\mathbf{Obs}_{sum}} },\frac{N_{\mathbf{Obs}_z} }{N_{\mathbf{Obs}_{sum}} }\right] \\
\mathbf{Cov}_{rot} &= 3\left[\frac{N_{\mathbf{Obs}_{roll}} }{N_{\mathbf{Obs}_{sum}} },\frac{N_{\mathbf{Obs}_{pitch}} }{N_{\mathbf{Obs}_{sum}} },\frac{N_{\mathbf{Obs}_{yaw}} }{N_{\mathbf{Obs}_{sum}} }\right]
\end{align}


Here, \( N_{\mathbf{Obs}_x}, N_{\mathbf{Obs}_y}, N_{\mathbf{Obs}_z}, N_{\mathbf{Obs}_{roll}}, N_{\mathbf{Obs}_{pitch}}, \) and \( N_{\mathbf{Obs}_{yaw}} \) represent the counts of observability labels for translational and rotational components, while \( N_{\mathbf{Obs}_{sum}} \) denotes the total number of labels across all state directions. These equations yield a normalized measure of observability for each component, enabling construction of a 6x6 diagonal covariance matrix \( \mathbf{Cov} \).


It is important to note that these equations evaluate the ratio of observability labels in each direction relative to the total number of labels. This relative framework operates on the hypothesis that, in well-structured environments, observability should be uniformly distributed across all state directions. Using this assumption, the method determines whether the current distribution of observability labels is balanced and quantifies its deviation from uniformity.

By examining the distribution of observability labels across the scan correspondences, this approach identifies under-constrained state directions that may require additional odometry to improve reliability. We eliminate the need to define arbitrary thresholds for observability across different methods, thereby enhancing the generality and robustness of the approach.

For example, the state direction along a long corridor is poorly constrained by LiDAR odometry. However, our proposed observability analysis of geometric features can predict degraded directions in advance. This allows the method to proactively incorporate pose prior regularization from an alternative odometry source specifically for ill-conditioned directions, rather than applying it indiscriminately to the entire state. This "active" fusion strategy has demonstrated greater effectiveness, significantly enhancing the robustness of SLAM systems.



%Taking point to plane distance as example, 
% We try to find a rigid-body transformation, composed of a rotation \( R \) and a translation \( \mathbf{t} \), that minimizes the sum of squared distances of each \( \mathbf{p}_i \) to the plane tangent to \( Q \) at \( \mathbf{q}_i \). The alignment error is given by:
% \begin{equation}
% E = \sum_{i=1}^{k} ((R\mathbf{p}_i + \mathbf{t} - \mathbf{q}_i) \cdot \mathbf{n}_i)^2
% \end{equation}
% Substituting and expanding, we can find the  error function \( E \) is :
% \begin{equation}
% E = \sum_{i=1}^{K} \left((\mathbf{p}_i - \mathbf{q}_i) \cdot \mathbf{n}_i + \mathbf{r} \cdot (\mathbf{p}_i \times \mathbf{n}_i) + \mathbf{t} \cdot \mathbf{n}_i\right)^2
% \label{eq:3}
% \end{equation}
% From Equation \ref{eq:3}, points with normal vectors \(\mathbf{n}_i\) are perpendicular to \(\mathbf{t}\), or torque vectors \(\mathbf{p}_i \times \mathbf{n}_i\) are perpendicular to \(\mathbf{r}\), do not change the error \(E\) (Figure 2 (a), (c)). More generally, \(\Delta d_i\) is zero for point-pairs whose 6-vector of torques and forces is orthogonal to the transformation vector. 

% We can calculate the observability of raw sensor measurements in the front end on each axis. For each correspondence point $\mathbf{p}_i$ in the source cloud $P$, we compute six values and project it to the world frame $[\mathbf{X}_w,\mathbf{Y}_w,\mathbf{Z}_w, \mathbf{roll}_w,\mathbf{pitch}_w, \mathbf{yaw}_w]$ as uncertainty for each direction.

% Taking point-to-plane distance as example, for every point we keep the planar scalar $a_{2D}$ of its neighborhood, as defined by \cite{deschaud2018imls}: $a_{2D} = (\sigma_2 - \sigma_3)/\sigma_1$ where $\sigma_i = \sqrt{\lambda_i}$ and $\lambda_i$ are eigenvalues of the PCA for the normal computation. Second, we compute the six values for every point $x_i$ to evaluate the observability of unknown rotation (first three) and translation (last three) shown in \eqref{eq:obs}.

\textbf{Adaptive Engine Selection~}
% Some methods perform observability analysis at the back-end level, which can be too late to predict potential failures \cite{degeneracy}, since insufficient data association already leads unreliable state estimation.
% By contrast, we perform observability analysis on the front-end level and evaluate the observability of feature detection, so that we can predict the degradation of LiDAR or visual odometry.
In real-world environments, sensors often operate in off-nominal conditions for extended periods, such as when a robot navigates through a dark environment. To reduce computational load and avoid outliers from malfunctioning sensors, we propose an adaptive engine (factor) selection approach to address long-term sensor-level degradation. We believe that factor graph optimization should be dynamically reconfigurable based on the type of sensor degradation, as shown in Fig.~\ref{fig:method} C-Step3. By evaluating the quality of each factor, the optimization process determines whether to include or exclude a given sensor’s data.

For example, in cases of prolonged visual degradation, the factor graph will \textit{reject all visual factors and rely solely on LiDAR}. Conversely, if LiDAR data becomes unreliable (i.e., all state directions are unobservable), the optimization will reject LiDAR factors. In mixed degradation scenarios, the factor graph will automatically select the most resilient sensor data. For complete sensor failure, the system will depend only on IMU factors, provided by learning IMU odometry.

We adopt a factor graph approach to account for all residual terms. However, not all residuals need to be included in the optimization, particularly in degeneracy-prone environments. Guided by the prior observability analysis, we utilize the covariance matrix $\mathbf{cov}$  to determine whether a specific factor should be incorporated into the optimization.

% \begin{equation}
% \begin{split}
% \mathop {\min }\limits_{{\mathbf{T}_{i + 1}}} \left\{
% \sum\limits_{p \in {\mathbb{F}_{i}}} {{{\left\| {e_i^{po \to po,li,pl}} \right\|}_{\mathbf{cov_L}}^2}} 
% + \sum\limits_{\left( {i,i + 1} \right) \in \mathcal{B}} {{{\left\| {{e^{imu}}} \right\|}_{\mathbf{cov_I}}^2}} \right. \\
% \left. + \sum\limits_{\left( {i,i + 1} \right) \in \mathcal{B}} {{{\left\| {{e_{io}^{prior}}} \right\|}_{\mathbf{cov_{i\_prior}}}^2}} 
% + \sum\limits_{\left( {i,i + 1} \right) \in \mathcal{B}} {{{\left\| {{e_{vo}^{prior}}} \right\|}_{\mathbf{cov_{v\_prior}}}^2}} 
% \right\}
% \end{split}
% \end{equation}


\begin{equation}
\begin{aligned}
\min_{\mathbf{T}_{i+1}} \bigg\{ 
& \sum_{p \in \mathbb{F}_i} \left\| e_i^{po \to po,li,pl} \right\|_{\mathbf{cov_L}}^2 
+ \sum_{(i, i+1) \in \mathcal{B}} \left\| e_{io}^{prior} \right\|_{\mathbf{cov_{io}}}^2 \\
& + \sum_{(i, i+1) \in \mathcal{B}} \left\| e_{vo}^{prior} \right\|_{\mathbf{cov_{vo}}}^2 
\bigg\}
\end{aligned}
\end{equation}







In this context, \( e^{prior}_{io} \) represents the relative pose factor provided by learning-based IMU odometry, weighted by the proposed \( \mathbf{cov_{io}} \). Similarly, \( e^{prior}_{vo} \) is the relative pose factor from visual odometry, weighted by \( \mathbf{cov_{vo}} \). The term \( e_i^{po \to po,li,pl} \) indicates point-to-point (line, distance) residuals between two keyframes with covariance \( \mathbf{cov_{L}} \)\cite{zhao2021super}. 


% Additionally, \( e^{imu} \) denotes the IMU preintegration factor with covariance \( \mathbf{cov_{I}} \)\cite{Forster_2015}. 
For further details on the \textbf{relative pose factor}, please refer to the Supplementary. This adaptability ensures that our sensor fusion pipeline remains flexible and resilient, dynamically reconfiguring itself to optimize performance under varying conditions.



\textbf{Adaptive Parameter Tuning~}
As shown in Fig.~\ref{fig:method} C-Step1, our method incorporates a series of self-tuning strategies to ensure adaptability across diverse platforms and environments \cite{ramezani2022wildcat} \cite{lim2023adalio}. For visual odometry, we adjust image contrast to enhance feature detection \cite{bergmann2017online}. However, LiDAR odometry can experience divergence in narrow, confined spaces, such as spiral stairs or corridors, due to fixed parameters that are typically optimized for open environments. To address this, our algorithm dynamically adjusts the weighting between visual and LiDAR odometry, enabling it to adapt to both narrow and open spaces. For a more detailed illustration, please refer to \textbf{Adaptive State Direction in Multi-floor Environments Experiment} in Supplementary.

Additionally, other parameters can be adjusted such as the size of voxelization, which can be adapted according to the scale of the environment. The algorithm also utilizes intensity values to filter airborne obscurants\cite{best2022resilient} and the gravity vector to estimate the initial rotational alignment between the LiDAR and IMU\cite{nemiroff2023joint}. These strategies significantly enhance the robustness of the system across a variety of platforms and environments. 




\subsection{Learning-based Inertial Odometry} 

We introduce a novel deep learning-based inertial odometry designed to capture and learn the motion patterns of various robots. The structure of the proposed system is illustrated in \fref{fig:method} A. The learning-based IMU odometry comprises two key components: an \textbf{pre-trained IMU model} trained from large-scale datasets, and \textbf{an adapter network} designed to handle previously unseen environments.

To achieve broad generalization and few-shot capabilities, our pre-trained IMU model is trained on a wide range of diverse datasets, including data from wheeled, legged, and aerial robots, as well as handheld devices. The model consists of five primary modules: IMU data preprocessing, data augmentation, a heterogeneous shared IMU backbone, a multihead layer design, and a loss function shown in \fref{fig:method} B. For more detail, please refer to our chapter 6 TartanIMU.



% \paragraph{Data Preprocessing}
% Our training dataset comprises over 100 hours of real-world IMU data, primarily sourced from our previous works: SubT-MRS \cite{zhao2024subt}, Tartan Drive \cite{sivaprakasam2024tartandrive}, IDOL \cite{sun2021idol}, Blackbird \cite{antonini2018blackbird}, and the UZH dataset \cite{Delmerico19icra}. This dataset encompasses eight distinct robot platforms with varying speeds and dynamics, including wheeled robots, drones, legged robots, and human-held data. It includes both raw IMU data and the ground truth trajectories of the robots.

% Given that these datasets use different IMU sensors with varying axis orientations, we first convert all axes to a consistent coordinate system, defined as follows: X points forward, Y points left, and Z points upward. Next, we standardize the IMU sampling frequency to 200 Hz across all datasets using interpolation. Finally, we filter out IMU data with significant noise by evaluating the IMU bias. This preprocessing is crucial for providing a consistent input frequency, reducing variability caused by differing sampling rates, and ensuring uniformity during training and inference. 

% \paragraph{Data Augmentation}
% Inspired by the work of Cao et al. \cite{cao2022rio}, we rotate the IMU data around the X, Y, and Z axes by a specified degree. Simultaneously, the corresponding ground truth trajectory is also rotated by the same angle. A pre-trained neural network for the IMU is trained on both the original and rotated datasets, as illustrated in Fig.~\ref{fig:method} B. This approach significantly enlarges the IMU training dataset and enhances the model's generalization capability.  

% \paragraph{Heterogeneous Shared Backbone}
% After preprocessing and augmentation, the IMU data is fed into the heterogeneous shared backbone. This backbone is designed to work across diverse robot platforms, creating a shared latent feature space. This space maps IMU data from various robots into a unified high-dimensional representation, enabling the model to capture "common knowledge" of robot motion patterns. By leveraging this shared representation, IMU data from new platforms can be integrated with minimal additional data and training, allowing the model to learn new motion patterns efficiently. Consequently, the heterogeneous backbone is capable of generalizing across a wide range of robots, capturing broader, world-level motion knowledge.

% As illustrated in \fref{fig:method} B, the network architecture consists of three main components: a 1D ResNet, standard LSTM layers, and a multi-head layer. The ResNet serves as the shared backbone, mapping IMU data from different robots into the unified high-dimensional space. It extracts latent features that represent the robot's motion. The LSTM then integrates the current latent features with the hidden state from the previous time step, capturing temporal dependencies in the motion data. Finally, the multi-head layers predict the local velocity for the input window, along with the corresponding covariance, completing the motion estimation process.

% \paragraph{Multi-head}
% Building upon the heterogeneous shared backbone, we propose a multi-head design to map the latent representation space to the velocity estimation space. Each head is designed to different robot types, benefiting from joint training within the shared backbone, as illustrated in \fref{fig:method} B. This multi-head architecture offers two key advantages: (i) it allows the model to learn diverse motion patterns in parallel, enhancing adaptability across various robots during deployment; and (ii) it stabilizes the training process by preventing competing learning objectives (e.g., drone motion is completely different from car motion.)

% Multi-head captures the distinct motion patterns of different robots, enabling parallel learning of diverse movements. This flexibility enhances the model's adaptability to various robots in real-world deployment. The mathematical formulation of the network is as follows:

% \begin{equation}
% \begin{aligned}
% (\hat{v}, \hat{u}) &= f\left(\left( ^B \mathbf{a}_{n-N}, ^B \boldsymbol{\omega}_{n-N} \right), \dots, \left( ^B \mathbf{a}_n, ^B \boldsymbol{\omega}_n \right), \mathbf{h}_{n-N} \right) \\
% ^B \mathbf{a}_n &= ^B_W\mathbf{R}_n (^W_B \mathbf{R}_n \left( \mathbf{a} - \mathbf{b}_a \right) - ^W \mathbf{g} )\\
% ^B \boldsymbol{\omega}_n &= \boldsymbol
% {\omega} - \mathbf{b}_g 
% \end{aligned}
% \end{equation}
 
% Here, \( f(\cdot) \) represents the function defined by the neural network that processes inputs from the IMU sensor. The network receives as input the acceleration \(\mathbf{a}\), angular velocity \(\boldsymbol{\omega}\), and the hidden state \(\mathbf{h}_{n-N}\) produced by the LSTM at the previous time step. For each IMU measurement, we remove the gravity vector \( ^W \mathbf{g} = [0, 0, 9.8] \) as well as the acceleration bias \(\mathbf{b}_a\) and gyro bias \(\mathbf{b}_g\). At each time step, the network predicts the current motion based on the hidden state \(\mathbf{h}_{n-N}\) and a local window of \(N\) samples of acceleration and angular velocity in the inertial frame \(B\). The output of the network consists of the estimated relative velocity \(\hat{{v}}\) and their associated uncertainties \(\hat{{u}}\).

% Unlike existing learning-based IMU odometry methods \cite{liu2020tlio}\cite{chen2021rnin}, which process IMU data in world coordinates \(W\). We process the raw IMU measurements in the \textbf{body frame} \(B\), leading to improved generalization capabilities.


 


% \paragraph{Loss Function}
% Our network employs relative loss functions, as illustrated in \fref{fig:method}~B. The relative loss function helps the network capture motion dynamics within a single window, improving the smoothness of the predicted odometry. To optimize the relative loss, we utilize the Mean Squared Error (MSE), defined as follows:
% \begin{equation}
% L_{RL}^{MSE}(\mathbf{v}, \hat{\mathbf{v}}) = \frac{1}{n} \sum_{i=1}^{n} \left(\sum_{j=m}^{m+M} \left( v_{j \rightarrow j+1} - \hat{v}_{j \rightarrow j+1} \right)\right) 
% \end{equation}

% In contrast to existing works \cite{liu2020tlio}\cite{chen2021rnin}, we define the loss in the \textbf{body frame} rather than in world coordinates, enabling the model to learn more generalized features. Specifically, \(\hat{v}_{j \rightarrow j+1}\) represents the 3D velocity output of the network for the \(j\)-th window in the body frame, while \(v_{j \rightarrow j+1}\) denotes the corresponding ground truth in the same frame. The variable \(n\) refers to the batch size during training, \(m\) is the starting window of the sequence, and \(M\) indicates the total number of LSTM windows.


% % \textit{Absolute Loss:} While achieving accuracy within individual windows is essential, minimizing cumulative errors over longer periods is equally critical. To enable the network to capture long-term motion patterns and mitigate error accumulation, we introduce an absolute loss function defined as:
% % \begin{equation}
% % L_{AL}^{MSE}(\mathbf{v}, \hat{\mathbf{v}}) = \frac{1}{n} \sum_{i=1}^{n} \left(\sum_{j=m}^{m+L} \left( v_{m \rightarrow j+1} - \hat{v}_{m \rightarrow j+1} \right)\right) 
% % \end{equation}

% \textit{Covariance:} The covariance estimation follows the methodology outlined in \cite{liu2020tlio}. The network outputs a three-dimensional vector, where each element represents the logarithm of a diagonal element of the covariance matrix $\Sigma $.

% \begin{equation}
% L^{\text{NLL}} = \frac{1}{2} (v - \hat{v})^T \Sigma^{-1} (v - \hat{v}) + \frac{1}{2} \ln \left| \Sigma \right|
% \end{equation}

% For more \textbf{implementation details} on learning imu odometry, please refer to our Chapter 6. 

% \paragraph{Adaptor Network}
% Although our pre-trained IMU model is designed to generalize across various platforms, it cannot guarantee 'few-shot' capability for every robot or motion pattern. To address this limitation, we introduce a module that enables the pre-trained IMU model to acquire new knowledge and continuously improve as it processes additional data. This adaptive capability is facilitated by a lightweight Low-Rank Adaptation (LoRA) adapter network, which is integrated into the pre-trained IMU model, as illustrated in Fig.~\ref{fig:method} B.

% To update each weight matrix \( W_0 \in \mathbb{R}^{d \times k} \) of the pre-trained IMU model, we apply a low-rank decomposition: 
% \begin{equation}
% W_0 + \Delta W = W_0 + BA,
% \end{equation}
% where \( B \in \mathbb{R}^{d \times r} \) and \( A \in \mathbb{R}^{r \times k} \), and $r$ is the rank and \( r \ll \min(d, k) \). During adaptation, \( W_0 \) remains fixed, while \( A \) and \( B \) are trainable parameters. Both \( W_0 \) and \( \Delta W = BA \) process the input, and their outputs are summed, resulting in:
% \begin{equation}
% h = W_0 x + \Delta W x = W_0 x + BA x    
% \end{equation}

% where \( A \) is initialized by a random Gaussian distribution and \( B \) starts at zero. %We then scale \( \Delta W x \) by \(\alpha \) to control the adaptation process, where $\alpha$ is a constant. 

% By freezing the parameters of the pre-trained model, our adaptive sensor fusion framework exclusively transfers knowledge to the LoRA adapter network \cite{hu2021lora}, which helps prevent catastrophic forgetting. The relatively small parameter set of the adapter network allows for the rapid assimilation of new information, enabling real-time adaptation. 
% For a more detailed illustration, please refer to our experiments on \textbf{catastrophic forgetting} in Supplementary Material.

% This approach ensures flexible adaptation to diverse robotic systems while preserving the strong generalization ability of the original pre-trained IMU model. It allows for efficient incorporation of new knowledge and system-specific optimization, without compromising the core capabilities of the pre-trained network.

\section{Evaluation}

\begin{figure*}[p]
   \centering
\includegraphics[width=0.95\textwidth, height=0.9\textheight,keepaspectratio]{figs/ch7/science_robotics/large_cmu_campus.pdf}
% \setlength{\parskip}{1pt}
% \captionsetup{font=footnotesize, format=plain, labelfont=bf, skip=-1pt, parskip=-1pt}
 % \captionsetup{font=footnotesize, skip=0.1 pt}
    \caption{ \textbf{Evaluation of 13 types of degradation in a single run.} The color-coded trajectory depicts our estimated odometry of a legged robot navigating through over 13 complex degradation scenarios. Despite these difficulties, the final endpoint drift was only \textbf{20 cm} over a total distance of 2,966 meters.}
    \label{fig:comprehensive_result}
\end{figure*}

\subsection{Comprehensive Degradation in One Run} 

\paragraph{Carnegie Mellon University Campus}
To thoroughly evaluate the robustness of our odometry system, we conducted a comprehensive degradation experiment on a legged robot in a single run. This experiment included visual, geometrical, mixed, and complete degradation scenarios. The test route, located on Carnegie Mellon University's campus, covered a distance of 2966 meters and took approximately 46 minutes to complete. The route incorporates more than \textbf{13} complex degradation scenarios, involving various combinations of challenging conditions. These scenarios include geometrical degradations such as long, textureless corridors (Fig~\ref{fig:comprehensive_result} D), narrow and steep multi-floor environments (Fig~\ref{fig:comprehensive_result} E, F), and transitions between indoor and outdoor environments (Fig~\ref{fig:comprehensive_result} L). Additionally, we introduced visual degradations (Fig~\ref{fig:comprehensive_result} G-N) such as smoke, white walls, low lighting, darkness, glass corridors, featureless surfaces, and lens flare. Mixed degradations, like steep staircases in darkness (Fig~\ref{fig:comprehensive_result} B) and expansive dark courts (Fig~\ref{fig:comprehensive_result} C), were also tested. Finally, we examined complete degradation, where dense smoke enveloped the robot (Fig~\ref{fig:comprehensive_result} A).
 
Our odometry system successfully completed real-time state estimation without any failures. The final endpoint drift was only 0.2 m over a total distance of 2966 meters, resulting in an exceptionally low drift rate of \textbf{0.006\%}. The complete mapping results are shown in Fig~\ref{fig:comprehensive_result}, highlighting high-fidelity views of the map under selected degraded conditions. The mapping results indicate the high accuracy and robustness of our pose estimates.  Notably, this performance was achieved without using any backend optimization techniques, such as loop closure. For more details on results, please refer to Supplementary \textbf{Movie S3}.


\begin{figure*}[htbp]
   \centering
\includegraphics[width=0.68\textwidth,height=1.0\textheight,keepaspectratio]{figs/ch7/science_robotics/switch_scheme_between_learning_so.pdf}
    \caption{\textbf{Robust Performance in a Smoke Scenario.} \textbf{(A)} Adaptive engine selection between learning-based IMU and LiDAR odometry in challenging environments, with purple and blue trajectories indicating their respective use. \textbf{(B)} Trajectory and map generated using learning-based IMU odometry in heavy smoke. \textbf{(C)} Self-supervised adaptation fine-tunes the IMU model from "SubT car" to "Tartan Drive." \textbf{(D)} Confidence-driven adaptive fusion combines IMU and LiDAR odometry outputs using soft fusion, avoiding hard switches.}
    \label{fig:offroad_smoke}
\end{figure*}


\subsection{Enhance SLAM Robustness in Smoke Environments}  We design smoke experiments to demonstrate how learning-based IMU odometry can improve the robustness of SLAM systems. It includes two processes: \textbf{self-supervised online adaptation} (explained in \textbf{Chapter 6}) and \textbf{learning-based IMU model deployment} in the sensor fusion pipeline.

\textit{Self-supervised Online Adaptation: }
% As previously discussed, the proposed learning-based IMU odometry may not perform optimally across all robotic platforms due to unique robot dynamics and IMU placement differences. 
% To address this variability, an online training strategy is necessary to quickly  adapt our pretrained IMU model to a platform-specific IMU for optimal performance. 
Fig.~\ref{fig:offroad_smoke}(C) illustrates the self-supervised online adaptation process used to fine-tune our pre-trained IMU model from the SubT vehicle to the TartanDrive vehicle. Initially, the TartanDrive vehicle operates in a smoke-free environment, where our factor graph pipeline continuously provides supervisory signals to refine the learning-based IMU odometry. Remarkably, this online training process required only 45 seconds of data collection to successfully adapt the pre-trained "SubT" IMU model into the "TartanDrive" model.

\textit{Learning-based Inertial Odometry Deployment: } 
After fine-tuning the IMU model, we integrated the learning-based IMU odometry into our adaptive sensor fusion pipeline. We evaluated its performance in two challenging environments: an extreme smoke area and an open flat area, as depicted in Fig.~\ref{fig:offroad_smoke}A. In these environments, LiDAR odometry becomes highly unreliable. When the vehicle enters the smoke or open flat areas, our adaptive pipeline automatically prioritizes the learning-based IMU odometry solution (purple trajectory). This is because the learning-based IMU odometry remains unaffected by these environmental conditions and continues to provide reliable state estimation. Conversely, when the vehicle exits these areas and conditions improve, the pipeline seamlessly transitions back to the LiDAR odometry solution (blue trajectory). This switching behavior is governed by the confidence level of the LiDAR odometry, as illustrated in Fig.~\ref{fig:offroad_smoke}D.

Importantly, this transition is not a "hard switch" but a soft fusion approach that blends the outputs of both LiDAR and IMU odometry across each axis. This strategy mirrors the adaptive engine selection mechanism used in previous mixed degradation scenarios.


To validate our method, we conducted figure-eight driving tests, repeatedly transitioning between smoky and clear areas. The adaptive switching occurred multiple times, demonstrating the pipeline's responsiveness to environmental changes. The estimated trajectory and the corresponding map are shown in Fig.~\ref{fig:offroad_smoke}B. While the reconstructed map exhibits significant noise due to the smoke, our sensor fusion pipeline consistently delivered robust state estimation.

This experiment demonstrates that incorporating learning-based IMU odometry significantly enhances the robustness of SLAM systems in severely degraded environments. For more details on results, please refer to Supplementary \textbf{Movie S6}.


\section{Conlusion}
In this section, the author provides insights and lessons learned on developing robust odometry systems in degenerate environments.

\paragraph{Hierarchical Adaptation is a Key Factor for Resilience} State estimation in challenging environments requires not only sensor redundancy but also computational efficiency. Sensor fusion, therefore, should be more intelligent and cannot be achieved by simply adding more sensors~\cite{ebadi2021dare}. However, much of the current SLAM literature focuses on "rigid" frameworks that prioritize robustness through the fusion of multiple modalities, rather than fostering resilience and online adaptability~\cite{prorok2021beyond}. This highlights a critical need for foundational research on resilient systems~\cite{ebadi2023present}.

To address this gap, we propose a hierarchical adaptation framework to build a resilient odometry system. The hierarchical structure is inspired by the way humans and animals have evolved to respond to varying levels of degradation. Given that environmental degradation can range from mild to extreme, we believe the odometry system should be dynamically reconfigurable, functioning as a multi-level mechanism. It should begin with low-level, resource-efficient adaptations and escalate to higher-level, computationally intensive adaptations as needed. When lower-level mechanisms are insufficient, higher-level cognitive processes intervene to ensure system recovery. Therefore, our approach dynamically selects the most relevant features, state directions, sensor engines, and deep networks—striking a balance between data redundancy for robustness and computational efficiency.

By incorporating multiple levels of adaptive strategies, the system ensures reliable state estimation across a wide range of challenging conditions, including visual, geometric, mixed, and complete sensor degradation (as illustrated in Fig.~\ref{fig:comprehensive_result}).

\paragraph{Low-cost IMUs are More Useful than Expected} Historically, IMUs usually have been regarded as auxiliary or secondary sensors within most SLAM pipelines, typically limited to pre-integration ~\cite{forster2015imu} to assist other sensors in the fusion process. However, we argue that IMUs hold far greater potential and can function as independent sensors in state estimation,  which is the same as LiDAR and visual sensors. 

 For the \textbf{first time}, our approach elevates the role of low-cost IMUs as primary navigation tools in severely degraded environments. We propose a heterogeneous, learning-based IMU odometry framework capable of providing generalizable odometry estimation across diverse robotic platforms. This framework integrates seamlessly into a hierarchical adaptation pipeline that supports online learning, enabling continuous adaptation to new and dynamic environments.
 
In scenarios where visual and LiDAR sensors are compromised—such as extreme conditions—the IMU odometry serves as the primary method for state estimation. This redundancy and fault tolerance significantly enhance the overall robustness of the system, ensuring reliable performance even under adverse circumstances.

\paragraph{Improving Sensor Calibration and Timesync} Future improvements should address several critical aspects, including sensor calibration and time synchronization \cite{lv2022observability,10443246}. Currently, the performance of our system depends heavily on accurate calibration and precise time synchronization. Implementing advanced online calibration and time-sync techniques could eliminate the need for manual parameter tuning. These methods would ensure that sensor data from different sources are accurately aligned in both time and space, thereby providing consistent and correct constraints for the sensor fusion process.

\paragraph{Improving Generalizability of the IMU model} A key area for improvement is the faster adaptation of the learning-based IMU odometry. While the current model has been evaluated across various robotic platforms, it may encounter challenges in quickly adapting to unseen robots and providing reliable state estimation in diverse environments. This limitation is primarily attributed to the domain gap between the training and testing data distributions. Leveraging IMU data from both real-world scenarios and simulations presents a promising direction for developing a more generalizable model, thereby reducing the discrepancy between training and testing phases.


In conclusion, the proposed hierarchical adaptation system shows significant potential for achieving robust state estimation in a variety of degraded environments. Despite remaining challenges—such as precise time synchronization, accurate calibration, and enhancing the generalization of the learning-based inertial odometry—this research lays a strong foundation for advancing resilient algorithms and systems for state estimation.
Overall, this work makes a valuable contribution to the growing field of adaptive odometry under challenging conditions. The system’s robustness, adaptability, and efficiency provide promising solutions for sustained and reliable operation in complex and degraded environments.
